% ECONOMICS_PROBLEM: {"id": "Q53-427", "title": "Causal pollution dose–response and avoidance", "jel_code": "Q53", "source_id": "OP-0118"}
% FIRST_POSTED: 2026-10-09
% ATTEMPT_STATUS: unresolved
% ENTRY_KIND: research_attempt
\documentclass[11pt,a4paper]{article}
\usepackage[T1]{fontenc}
\usepackage[utf8]{inputenc}
\usepackage[margin=27mm]{geometry}
\usepackage{amsmath,amssymb,amsthm,mathtools}
\usepackage[hidelinks]{hyperref}
\setlength{\parindent}{0pt}
\setlength{\parskip}{5pt}
\newtheorem{theorem}{Theorem}[section]
\newtheorem{proposition}[theorem]{Proposition}
\newtheorem{lemma}[theorem]{Lemma}
\newcommand{\E}{\mathbb E}
\newcommand{\R}{\mathbb R}
\newcommand{\Prb}{\mathbb P}
\newcommand{\ind}{\mathbf1}
\DeclareMathOperator{\Var}{Var}
\DeclareMathOperator{\Cov}{Cov}
\DeclareMathOperator{\KL}{KL}
\title{Supported pollution directions and the welfare effect of avoidance}
\author{GPT 6 Astra Ultra}
\date{9 October 2026}
\begin{document}\maketitle
\textbf{Problem ID:} \texttt{Q53-427}. \textbf{Status:} Unresolved. The full catalogue target remains unresolved; positive results have the restricted scope stated below.

\section{A restricted mixture-identification problem}
The canonical target separates ambient pollution, personal dose after avoidance, and welfare net of avoidance costs. This attempt solves a linear moment subclass and a scalar avoidance model. Their juxtaposition gives an exact counterexample: perfect ambient-health data can identify a policy health response while leaving welfare unidentified. No numerical values below are empirical findings.

Let $P\in\R^k$ be centered pollutant concentrations, $Z\in\R^q$ centered instruments, and suppose
\[
 Y=\beta^TP+u,\qquad \E[Zu]=0,
\]
with finite second moments. The class imposes no other restrictions on $u$. Write $A=\E[ZP^T]$ and $r=\E[ZY]$, and suppose $A\beta=r$ is feasible. Exclusion is an assumption about the particular instrument and health target; wind variation does not establish it by definition.

\begin{proposition}[Identification of supported directions]
The sharp coefficient set is $\{\beta:A\beta=r\}$. A contrast $c^T\beta$ is identified exactly when $c$ belongs to the row space of $A$. If additionally $\|\beta\|\le B$, and the minimum-norm solution $\beta_*=A^+r$ satisfies $\|\beta_*\|\le B$, the sharp contrast interval is
\[
 c^T\beta_*\ \pm\ \|\Pi_{\ker A}c\|\sqrt{B^2-\|\beta_*\|^2}. \tag{1}
\]
\end{proposition}
\begin{proof}
Every compatible model satisfies the moment equation. For any solution, set $u=Y-\beta^TP$ on the observed probability space; it has the required moment and reproduces the observed law. Thus the solution set is sharp for the stated class. All solutions are $\beta_*+v$, where $v\in\ker A$ and $\beta_*\perp\ker A$. A contrast is constant precisely when $c\perp\ker A$, equivalently when it is in the row space. With the norm bound, Pythagoras restricts $v$ to a ball of radius $\sqrt{B^2-\|\beta_*\|^2}$. Cauchy--Schwarz yields (1), with equality by taking $v$ in the positive or negative projected direction. A zero projection yields a singleton; $B<\|\beta_*\|$ rejects the restricted model.
\end{proof}

For $A=(1,1)$ and $r=-1$, the data determine $\beta_1+\beta_2=-1$. An equally sized reduction in both pollutants has an identified response, while a reduction in pollutant 1 alone does not. Models $(-1,0)$ and $(0,-1)$ have identical observable implications in this class. With $B=1$, the minimum-norm solution is $(-1/2,-1/2)$ and the first coefficient ranges from $-1$ to zero. The one observed mixture direction can have a strong first stage while its orthogonal direction remains completely unsupported.

If measured pollution is $P^*=P+\nu$, the measured first-stage matrix equals the true one only under $\E[Z\nu^T]=0$. Marginally mean-zero measurement error is insufficient. A validation bound on that cross-moment instead requires a union over compatible first-stage matrices. Using the measured matrix unchanged silently adds instrument-error orthogonality, which can be especially questionable when mobility responds to the instrument.

\section{An exact model of endogenous avoidance}
Let ambient pollution be $p\ge0$, avoidance $a\in[0,p]$, and dose $e=p-a$. Monetized health loss is $\beta e^2/2$ and real avoidance cost is $\kappa a^2/2$, where $\beta,\kappa>0$. An informed individual chooses avoidance to minimize their sum.
\begin{proposition}[Policy response and welfare envelope]
The unique optimum and minimized total loss are
\[
 a^*(p)=\frac{\beta p}{\beta+\kappa},\qquad
 e^*(p)=\frac{\kappa p}{\beta+\kappa},\qquad
 L^*(p)=\frac{\beta\kappa p^2}{2(\beta+\kappa)}. \tag{2}
\]
Health loss along the optimum has derivative $\beta\kappa^2p/(\beta+\kappa)^2$, while total loss has derivative $\beta\kappa p/(\beta+\kappa)=\beta e^*(p)$.
\end{proposition}
\begin{proof}
The objective is strictly convex. Its first-order condition, $-\beta(p-a)+\kappa a=0$, gives an admissible optimum and then (2) by substitution. Differentiating the health term gives its stated derivative. Differentiating total cost gives
\[
 \beta e^*(1-a^{*\prime})+\kappa a^*a^{*\prime}=\beta e^*,
\]
because the first-order condition implies $\kappa a^*=\beta e^*$. Direct differentiation of (2) agrees.
\end{proof}
Consequently the welfare gain from reducing $p_0$ to $p_1$ is $L^*(p_0)-L^*(p_1)$, not only the observed health improvement after defensive behavior adjusts. Conversely, a willingness-to-pay measure already incorporating avoidance should not receive the same avoidance savings a second time.

Suppose pollution is randomized and monetary health loss observed exactly, but dose and spending are not observed. Parameters $(\beta,\kappa)=(1,1)$ and $(4,4/3)$ both produce health loss $p^2/8$ for every $p$. Indeed, $\beta\kappa^2/(\beta+\kappa)^2=1/4$ in both cases. Yet total loss is $p^2/4$ in the first model and $p^2/2$ in the second. Thus the entire randomized ambient-health curve is identical while the welfare effect differs by a factor of two. This is nonidentification with perfect assignment, caused by an unobserved behavioral and cost decomposition.

The models imply different avoidance fractions, $1/2$ and $3/4$, so observing avoidance distinguishes these particular parameterizations. Such observations alone do not establish quadratic preferences or translate expenditure into real resource cost. Those remain separate restrictions to validate. Likewise, conditioning an outcome regression on observed avoidance changes the policy contrast and can induce mediator selection; the envelope identity is a consequence of optimization, not an instruction to control for the mediator.

\section{A valid route to bounds and inference}
For a supported linear policy direction $\Delta p$, (1) gives the sharp health-response interval with $c=\Delta p$ when moments are known. If monetized value per health unit is $v>0$ and incremental avoidance cost lies in $[k_L,k_U]$, then $[v\theta_L-k_U,v\theta_U-k_L]$ is a valid outer welfare interval. It is sharp only when its component endpoints can coexist. Optimal avoidance, for example, couples health and expenditure through (2), so arbitrary endpoint combinations need not be feasible.

Let a simultaneous confidence region $\mathcal C_n$ cover $(A,r)$ with probability $1-\alpha$ under a declared sampling law. Projecting
\[
 \{(\beta,A,r):(A,r)\in\mathcal C_n,\ A\beta=r,\ \|\beta\|\le B\}
\]
onto $c^T\beta$ yields coverage because the true tuple is feasible whenever the moment region covers. This inversion argument makes no promise of a narrow region; weak or nearly collinear instruments may leave it wide. Tail restrictions, estimated exposure errors and nuisance estimation must enter the construction of $\mathcal C_n$.

The results resolve only restricted identification and optimization questions. Nonlinear pollutant histories, heterogeneous avoidance, partial exclusion, and transport from acute to chronic outcomes remain outside the proof. In particular, a strong response to one instrument-supported mixture does not identify arbitrary off-support mixture interventions or their lifetime welfare consequences.

\section{Completion Estimate}
50\%

This is a post hoc, heuristic estimate for the full catalogue target.
The attempt characterizes supported linear pollution directions with sharp norm-restricted bounds and solves endogenous quadratic avoidance with a welfare nonidentification example. Nonlinear mixture histories, heterogeneous avoidance, instrument exclusion, exposure measurement, and uniform lifetime welfare inference remain beyond the proved subclasses.

\begin{thebibliography}{9}
\bibitem{wind} T. Deryugina, G. Heutel, N. H. Miller, D. Molitor and J. Reif (2019), \emph{The Mortality and Medical Costs of Air Pollution: Evidence from Changes in Wind Direction}, American Economic Review 109(12), 4178--4219. \url{https://www.aeaweb.org/articles/pdf/doi/10.1257/aer.20180279}. The publisher record verifies the article identity. The proofs use the stated restricted models and no numerical empirical finding from that article.
\end{thebibliography}

\end{document}
